% 3.1 — Besoins
\section{Analyse des besoins}
\label{sec:besoins}

\subsection{Acteurs}
Trois acteurs interagissent avec \demo{} : l'\textbf{agent}, qui propose des commandes et
des modifications de fichiers ; la \textbf{personne} (le développeur), qui installe l'outil,
répond aux questions qu'il pose et restaure ce qu'il a capturé ; et les \textbf{systèmes
tiers} touchés par les commandes (système de fichiers local, bases de données, dépôts git
distants, fournisseurs cloud).

\subsection{Besoins fonctionnels}
Le tableau~\ref{tab:bf} rassemble ce que l'outil doit faire pour ces acteurs.

\begin{table}[htbp]
\centering
\caption{Besoins fonctionnels}
\label{tab:bf}
\small
\begin{tabularx}{\linewidth}{@{}L{1cm}Y@{}}
\toprule
\textbf{Id} & \textbf{Besoin} \\
\midrule
BF1 & Intercepter, avant exécution, chaque commande shell (POSIX et PowerShell) et chaque
      modification de fichier proposée par l'agent, par toutes les portes disponibles. \\
BF2 & Prédire l'effet de la commande : ce qu'elle détruit, crée, déplace, et ce qu'elle
      atteint à distance. \\
BF3 & Rendre une décision parmi six : ALLOW, REVERSIBLE, DRY\_RUN, SUPERVISED,
      CONTEXT\_MISMATCH, ESCALATE ; pour une escalade, indiquer si une personne peut
      l'approuver sur le champ. \\
BF4 & Capturer, avant l'exécution, une copie exacte de ce qui sera détruit, en un seul
      groupe par commande, restaurable par \cmd{demo\_cli undo <id>}. \\
BF5 & Prévisualiser une requête SQL destructrice (nombre de lignes touchées). \\
BF6 & Placer sous un filet d'exécution ce qui ne peut pas être lu (Linux et Windows). \\
BF7 & Refuser toute modification des fichiers de \demo{} lui-même par l'agent. \\
BF8 & Enregistrer chaque décision dans un registre infalsifiable, vérifiable par
      \cmd{demo\_cli verify}. \\
BF9 & Diagnostiquer l'installation (\cmd{demo\_cli doctor}) et dire, par agent et par
      système, ce qui est protégé et ce qui ne l'est pas. \\
\bottomrule
\end{tabularx}
\end{table}

\subsection{Besoins non fonctionnels}
Le tableau~\ref{tab:bnf} fixe les qualités attendues, chacune avec sa justification.

\begin{table}[htbp]
\centering
\caption{Besoins non fonctionnels}
\label{tab:bnf}
\small
\begin{tabularx}{\linewidth}{@{}L{1.2cm}L{3.4cm}Y@{}}
\toprule
\textbf{Id} & \textbf{Qualité} & \textbf{Exigence et justification} \\
\midrule
BNF1 & Honnêteté & Une décision ne revendique jamais une protection non prise ; la raison
  dit ce qui a été capturé et ce qui ne l'a pas été. \\
BNF2 & Sûreté par défaut & Une \emph{décision} incertaine va vers le plus strict ; une
  ligne non lue n'est jamais ALLOW. \\
BNF3 & Ne pas bloquer l'agent sur une panne & Une erreur interne de \demo{} ne doit pas
  paralyser l'agent : elle laisse passer, bruyamment ; mais une panne du \emph{moteur} ne
  rend jamais une décision plus permissive (B0). \\
BNF4 & Latence & Le moteur dispose de 10\,ms de temps CPU (50\,ms pour une seconde
  tentative) et d'un plafond de 250\,ms en temps réel ; les familles d'outils ne se chargent
  qu'à la demande. Un outil lent est un outil désinstallé. \\
BNF5 & Portabilité & Linux et Windows (PowerShell, Git Bash), avec les mêmes garanties ;
  chaque partie est vérifiée sur les deux systèmes. \\
BNF6 & Non-régression & Aucune décision ne devient plus permissive sans arbitrage
  explicite (portail de \nbcorpus{} commandes). \\
BNF7 & Confidentialité & Les reçus masquent les secrets (mots de passe d'URL, clés) ;
  aucune télémétrie. \\
\bottomrule
\end{tabularx}
\end{table}
